\documentclass[12pt,a4paper]{article}
\usepackage[scheme=plain]{ctex}
\usepackage[margin=2.5cm]{geometry}
\usepackage{indentfirst}
\usepackage{setspace}
\usepackage{fontspec}
\setCJKmainfont[BoldFont=SimHei,ItalicFont=KaiTi]{SimSun}
\newCJKfontfamily\kaifont{KaiTi}[AutoFakeBold]
\usepackage{newpxtext,newpxmath}

\usepackage{lastpage}\usepackage{fancyhdr}





\pagestyle{plain}

\usepackage[nobottomtitles*]{titlesec}
\titleformat{\section}{\Large\bfseries}{}{0em}{ }
\titlespacing*{\section}{0pt}{3ex}{1.5ex}

\usepackage{xcolor}
\usepackage{needspace}

\setlength{\parindent}{2em}
\setlength{\parskip}{0.4em}

\doublespacing

\title{\bfseries Day 5\\[0.3em]教育作为一种行政现象}
\author{}
\date{}

\begin{document}
\maketitle

\vspace{-2em}

\vspace{1em}

\section{高考改卷组织与确定性偏好}
\vspace{-0.3em}{\color{gray}\small\kaifont \noindent 你们即将读到两种材料。一种是制度文件。《国家教育统一考试网上评卷工作考务管理办法》，教育部办公厅发布的。它规定了高考改卷的组织架构、人员选聘资格、评分细则制定流程、质量控制机制。它是一套精密机器的技术说明书。另一种是阅卷教师写的经验分享。它告诉你在这套机器里面工作是什么感觉。

\noindent 制度设计者说“这把尺子量到底是为了公平”，身处机器中的人说“我根本没时间欣赏你的文采”。两段话拼在一起，你看到的不再是公平或残酷的单一结论，而是一个组织行为问题：当“减少误差”成为唯一的优化目标，这套机制鼓励的不是最优答案，而是最可预测的答案。}\vspace{0.5em}
\vspace{0.3em}

\noindent\rule{\textwidth}{0.4pt}

\vspace{0.3em}

{\color{gray}\small （教育部办公厅，2006，2006年国家教育统一考试网上评卷工作考务管理办法）}



第三章　评卷工作的组织管理



第七条　评卷工作由省级教育考试机构负责组织和管理，有关高等学校给予支持配合。



第八条　评卷工作由省级教育考试机构在指定的高等学校或指定的其它场所集中进行。



第九条　评卷场地应有专人保卫，评卷员和评卷工作人员出入须凭《评卷工作证》，不得随意变更，也不得将《评卷工作证》转借他人。



评卷员使用的密码只限本人掌握，不得告知其他评卷员。



第十条　应成立由省级教育考试机构、评卷高校和学科负责人参加的评卷领导小组。



各学科（课程）成立评卷小组，负责人由本学科（课程）具有副教授以上职称的教师担任。



评卷程序和复评数量、标准等由评卷领导小组确定。



第十一条　评卷员由评卷领导小组聘任。评卷员应是从事本学科（课程）教学工作、业务水平较高、责任心强、作风正派的教师。评阅作文、论述等主观性试题的教师，必须具有本学科中、高级以上职称。



第十二条　评卷前应对全体评卷员及有关工作人员进行保密纪律的教育和评卷规定、评卷程序的培训。如果省级教育考试机构直接在社会上聘用评卷员，应与所聘任的评卷员签订评卷协议。



第十三条　评卷前各学科（课程）评卷小组负责人应组织本学科（课程）评卷员认真研究试题、答案和评分参考，并进行试评。在试评基础上制订评分细则，评分细则应报省级教育考试机构备案。



第十四条　评卷员必须严格执行评分细则，不得擅自更改。凡在评卷过程中对评分细则有异议，属于全国统一命题的科目的，须由省级教育考试机构报教育部考试中心（自考办）同意后方可改动；属于省自主命题的科目的，须报省级教育考试机构同意后方可改动。



第十五条　学科评卷小组负责人应熟悉评卷系统。评卷过程中，学科评卷小组负责人应随时掌握评卷进度和评卷员的评卷质量。对于不能准确把握评分细则的评卷员，经提醒后仍不能纠正的应及时进行调整。



第十六条　省级教育考试机构应与为网上评卷提供技术支持和服务的有关公司及所聘任的评卷工作人员签订评卷工作协议。









{\color{gray}\small （语文帮高中版，2026，高考阅卷不看文采、不看感悟，只看踩点给分（新课标一卷通用）}，今日头条）



一、先懂底层逻辑：高考阅卷，是标准化流水线，不是文学鉴赏





高考是国家级选拔性标准化考试，第一原则是：绝对公平、统一标准、可量化打分。





无数学生搞错了阅卷本质：



你以为阅卷老师是语文名师，慢慢品读你的答案、共情你的感悟、欣赏你的文笔；



真实阅卷现状： 每位阅卷老师日均阅卷3000–5000份，单道主观题阅卷时长3–8秒。





在极致的阅卷速度、统一的评分规则下：





1. 优美的辞藻、流畅的段落、细腻的个人感悟，无法量化，一律不计分；



2. 阅卷老师没有时间品读情感、赏析文笔，只做一件事：扫描关键词、核对得分点；



3. 所有主观题分数，全部绑定固定答题角度、专属得分术语、文本对应依据，缺一不得分。





简单一句话：高考语文的分数，是按点结算的，不是按情分、文笔分的。





你的个人感悟再深刻、文字再优美，只要没答出官方规定的得分点，就是0分；



你的答案朴实直白、毫无文采，只要精准踩中所有关键词、答题维度，就是满分。



{\color{gray}\small （天意，2026，高考阅卷老师的压力总结，小红书）}



额外的淘汰机制，全程都不能松懈

上岗前必须通过试评考核，打分和专家组标准误差超标，根本没有正式阅卷资格；阅卷中途一旦稳定性不达标，会被中途清退，这也让老师全程不敢放松心态。

简单总结：阅卷不是简单改卷子，相当于在严密的监控下，连续一周多做高精密的精细活，既要快、又要准，还要严守纪律，心理和身体都会达到负荷上限。



{\color{gray}\small （冰川思想库，2019，那些年，我经历过的高考阅卷，知乎）}



高考阅卷更算不得神圣。为什么参加高考阅卷？大多数人就是为了挣钱。





高考阅卷一般要连续三到四天，每天从早到晚，是比较辛苦的。老教师一般不愿参加，就是觉得累，挣钱还不多。那么，能挣多少钱呢？当时（十几年前）一次高考阅卷大概能拿到一千多元。





因为缺人，每次都有一定数量的研究生参与进来，对一些经济条件不好的研究生来说，几天时间能挣个千儿八百的还不错，据说，有的研究生还会特意争取阅卷的机会。





当然，没有高考阅卷经验的研究生的阅卷质量往往被重点关注，反复核查。具体到某个人高考阅卷的收入，主要是由产量（阅卷量）决定的。





每个人阅卷的数量都由专人统计。眼疾手快的人，可就占便宜了，阅卷收入要远远高于动作迟缓的人。

\newpage\section{街头官僚的职责拮抗}
\vspace{-0.3em}{\color{gray}\small\kaifont \noindent 前面我们看了高考改卷组织如何令可能怀有良好意愿的教师化为追求确定性的给分机器。你可能会问：那老师教课的时候呢？老师关心学生的时候呢？那些不在改卷体系管辖范围内的日常互动，难道也被"标准化"了吗？

\noindent 答案是没有。而这恰恰是问题所在。

\noindent 李普斯基在上世纪八十年代提出了"街头官僚"这个概念。他说的不是街上的公务员。他说的是那些在日常工作中直接面对公民、并对他人生活产生实际影响的一线工作者——警察、社工、法官、教师。这些人的核心特征不是职位高低，而是他们在执行政策的过程中拥有"不可消除的自由裁量权"——因为政策不可能覆盖人与人的每一次具体互动。

\noindent 这种自由裁量权带来的后果是什么？Lipsky 说，街头官僚同时承担着中心职责和边缘职责。中心职责是组织给他的任务，被上级严密管理、频繁考核。边缘职责则分散在日常互动的间隙中——它们不容易被识别、量化和管理，但同样真实地影响着他人的生活质量。这两类职责之间存在拮抗关系：管好了一个，另一个可能就会被压掉。}\vspace{0.5em}
\vspace{0.3em}

\noindent\rule{\textwidth}{0.4pt}

\vspace{0.3em}





街头官僚工作环境的一个根本特征是：他们必须面对服务对象对自己决定的个人反应，无论他们如何处理这些决定带来的后果。说人的自我评价受到街头官僚行为的影响，就等于说人对政策是反应性的。这不仅限于潜意识过程。街头官僚机构的服务对象会对真实的或感知到的不公正感到愤怒，会制定策略来讨好工作人员，会因街头官僚的决定而表现出感激和兴高采烈，或者闷闷不乐和消极被动。被电话公司、车管局或其他政府机构以一种冷漠和常规的方式对待是一回事——这些机构的代理人对你的具体情况一无所知。而被一个你正在与之面对面交谈、并期望至少得到一次开放而富有同情心的倾听的人以“官僚式的”（在这个词的贬义意义上）方式进行筛选、分类和对待，则完全是另一回事。简而言之，街头官僚的工作现实与官僚制所崇尚的非人格化超脱决策理想相去甚远。恰恰相反，在街头官僚机构中，关键决定的对象——人——实际上会因为这些决定而发生改变。



街头官僚也是公民反应的焦点，因为他们的自由裁量权开启了一种可能性：他们可能替某人做出有利的回应。他们对“公共利益”所承担的宽泛而模糊的义务，让人们怀抱希望，期待这个具体的工作人员会对自己的处境抱有好意或偏向。因此，在一个由庞大而非人格化的机构组成的世界中——这些机构表面上掌握着通往重要福利、制裁和机会的钥匙——工作定义的模糊性维系着“在这里或许能找到一个肯帮忙的人”的希望。
\par\vspace{0.3em}{\footnotesize\color{gray}（Lipsky, M. (1980). Street-Level Bureaucracy: Dilemmas of the Individual in Public Services. Russell Sage Foundation. 由 DeepSeek V4 Pro 翻译。）}
\newpage\section{教育作为排他性的国家治理领域}



\vspace{-0.3em}{\color{gray}\small\kaifont \noindent 1985年5月27日，中共中央发布了《关于教育体制改革的决定》。这是一个里程碑式的纲领性文件。它开篇就说："教育必须为社会主义建设服务，社会主义建设必须依靠教育。"它把教育定位为经济建设的人才供应线，要求培养"数以亿计"的劳动者、"数以千万计"的技术人员和各方面的工作者。教育行政组织与生活在学生中的教师的复杂心绪，从那一刻起，就被嵌入了一个更大的国家目标：快速实现现代化的经济和行政建制。}\vspace{0.5em}



\vspace{0.3em}

\noindent\rule{\textwidth}{0.4pt}

\vspace{0.3em}





{\color{gray}\small （中共中央，1985，中共中央关于教育体制改革的决定）}



党的十二届三中全会关于经济体制改革的决定，为我国社会生产力的大发展、为我国社会主义物质文明和精神文明的大提高，开辟了广阔的道路。今后事情成败的一个重要关键在于人才，而要解决人才问题，就必须使教育事业在经济发展的基础上有一个大的发展。

教育必须为社会主义建设服务，社会主义建设必须依靠教育。社会主义现代化建设的宏伟任务，要求我们不但必须放手使用和努力提高现有的人才，而且必须极大地提高全党对教育工作的认识，面向现代化、面向世界、面向未来，为九十年代以至下世纪初叶我国经济和社会的发展，大规模地准备新的能够坚持社会主义方向的各级各类合格人才。要造就数以亿计的工业、农业、商业等各行各业有文化、懂技术、业务熟练的劳动者。要造就数以千万计的具有现代科学技术和经营管理知识，具有开拓能力的厂长、经理、工程师、农艺师、经济师、会计师、统计师和其他经济、技术工作人员。还要造就数以千万计的能够适应现代科学文化发展和新技术革命要求的教育工作者、科学工作者、医务工作者、理论工作者、文化工作者、新闻和编辑出版工作者、法律工作者、外事工作者、军事工作者和各方面党政工作者。所有这些人才，都应该有理想、有道德、有文化、有纪律，热爱社会主义祖国和社会主义事业，具有为国家富强和人民富裕而艰苦奋斗的献身精神，都应该不断追求新知，具有实事求是、独立思考、勇于创造的科学精神。这就向我国教育事业的发展和教育体制的改革，提出了伟大而又艰巨的任务。

建国以来，我国教育事业的发展走过了曲折的道路。经过解放初期的接管改造和以高等学校院系调整为中心的教育改革，我们把旧中国的半殖民地半封建教育事业转变成为社会主义教育事业。三十几年来，依靠广大教育工作者的辛勤努力，教育事业取得了中国历史上从来没有过的巨大的发展，成绩是显著的。今天战斗在我国各条战线上的广大有文化的劳动者和各方面工作的骨干力量，绝大部分都是建国以后培养出来的。但是，另一方面，从五十年代后期开始，由于全党工作重点一直没有转移到经济建设上来，由于“阶级斗争为纲”的“左”的思想的影响，教育事业不但长期没有放到应有的重要地位，而且受到“左”的政治运动的频繁冲击。“文化大革命”更使这种“左”的错误走到否定知识、取消教育的极端，从而使教育事业遭到严重破坏，广大教育工作者遭受严重摧残，耽误了整整一代青少年的成长，并且使我国教育事业同世界发达国家之间在许多方面本来已经缩小的差距又拉大起来。

\end{document}
